\section{Introduction}
Une galerie classique retrouve facilement une photographie lorsque son nom ou sa catégorie est connu. En l'absence de cette information, une image d'exemple peut servir de requête. La recherche d'images par le contenu compare alors des caractéristiques calculées sur les pixels. Le classement dépend des caractéristiques choisies et peut privilégier des ressemblances différentes de celles que l'utilisateur juge pertinentes.

\projet{} répond aux deux tâches du mini-projet. La première consiste à gérer une galerie authentifiée, organiser les images en catégories, conserver les originaux lors des transformations et calculer les descripteurs de couleur et de texture. La seconde consiste à retrouver les images proches d'une requête et à affiner le classement lorsque l'utilisateur signale des exemples pertinents. La méthode affectée au \groupe{} est Re-weighting Type-1 \cite{renuka2016}.
